\documentclass[10pt,a4paper]{amsart}
\usepackage[margin=19mm]{geometry}
\usepackage{amsmath,amssymb,mathtools}
\usepackage[scr=rsfs,cal=euler]{mathalfa}
\usepackage{xfrac}
\usepackage[unicode,hidelinks,bookmarks=false]{hyperref}
\newcommand{\ie}{i.e.\ }
\newcommand{\cf}{cf.\ }
\newcommand{\resp}{respectively\ }
\newcommand{\Subsection}{\subsection}
\providecommand{\nobreakdash}{\nobreak\hbox{-}}
\setlength{\emergencystretch}{2em}
\multlinegap0pt
\usepackage{bm}
\usepackage{bbm}
\usepackage{dsfont}
\usepackage{xspace}
\usepackage{cancel}
\usepackage{mathtools}
\usepackage[shortlabels]{enumitem}
\usepackage{amsthm}
\makeatletter
\@ifundefined{Theorem}{\newtheorem{Theorem}{Theorem}}{}
\@ifundefined{Lemma}{\newtheorem{Lemma}[Theorem]{Lemma}}{}
\makeatother
\newcommand\CC{\mathbb{C}}
\newcommand\DD{\mathbb{D}}
\newcommand\ZZ{\mathbb{Z}}
\newcommand\eps{\varepsilon}
\newcommand\half{\ensuremath{\frac{1}{2}}}
\newcommand\numbers{\textup{(\arabic*)}}
\newcommand\thalf{\ensuremath{\tfrac12}}
\title[Positive Traces on Certain ${\rm SL}(2)$ Coulomb Branches]{Positive Traces on Certain ${\rm SL}(2)$ Coulomb Branches}
\author{Daniil Klyuev, Joseph Vulakh}
\date{Source: arXiv:2507.19447v2 (CC BY 4.0)}
\begin{document}
\maketitle

\noindent\emph{Source and licence.} Positive Traces on Certain ${\rm SL}(2)$ Coulomb Branches. Version arXiv:2507.19447v2 (CC BY 4.0), distributed under the Creative Commons Attribution 4.0 International licence, as recorded in the arXiv licence metadata for this exact version and shown on the abstract page https://arxiv.org/abs/2507.19447v2. Full rights notice: licenses/arxiv-2507.19447-CC-BY-4.0.txt.\par\smallskip

\noindent\textit{Editorial scope.} Counted unit: the Theorem whose source label is \texttt{thm:classif} in the article (source0.tex lines 719--752, source label \texttt{thm:classif}), delivered together with the article's own complete original proof of it (source0.tex lines 754--980, one \texttt{proof} environment of 227 lines). No mathematics is written, repaired or re-derived: statement and proof are the article's own text line for line. Required context retained uncounted: lem:a\_prop (lines 533--546); lem:xy\_poly (lines 588--612); lem:trace\_eh\_zero (lines 683--686). Every definition, display and label of those lines is retained unchanged. Omitted from this package and not redistributed: 1--532, 547--587, 613--682, 687--718, 753--753, 981--2462. Nothing inside a retained excerpt is omitted or altered: every retained body is delivered exactly as the article's lines are written, so no omission receipt of any kind accompanies this package. Citations: the retained excerpts carry no citation command at all, so no bibliography entry of the article is transcribed and no reference is redistributed. Reference adaptation: none was needed and none was made; every reference inside the delivered unit resolves inside it, so no unresolved reference or citation is produced. Printed slips kept literal: no printed word, symbol, hypothesis or number of the counted statement or of its proof was altered. Layout and class: the article's own class is \texttt{sigma} with packages including ifpdf, amsfonts, lipsum, enumitem, environ; the wrapper is the lane's pinned \texttt{amsart} editorial wrapper at 10pt on A4, which restates the theorem-like environments the retained text needs and adds the article's own macro definitions that the retained text needs (\texttt{\textbackslash CC}, \texttt{\textbackslash DD}, \texttt{\textbackslash ZZ}, \texttt{\textbackslash eps}, \texttt{\textbackslash half}, \texttt{\textbackslash numbers}, \texttt{\textbackslash thalf}), with extra wrapper packages (bm, bbm, dsfont, xspace, cancel, mathtools, enumitem). The delivered numbering is the unit's own. Assets: no retained span contains a figure environment, a graphics-inclusion command, a PDF-page inclusion, a table or a TikZ picture; no image file is copied into this package, the package declares no asset dependency and no image file is redistributed with it. Licence: version arXiv:2507.19447v2 of this article is distributed under the Creative Commons Attribution 4.0 International licence, as recorded in the arXiv licence metadata for this exact version and shown on the abstract page https://arxiv.org/abs/2507.19447v2. A full rights notice is delivered at licenses/arxiv-2507.19447-CC-BY-4.0.txt. Characters: every retained excerpt is pure ASCII, so the delivered engine, XeLaTeX with its default Latin Modern text font, reports no missing character.
\begin{Lemma}
 \label{lem:a_prop}
 The following statements hold in $H_c$:
 \begin{itemize}\itemsep=0pt
 \item $\alpha_q = -\alpha_{-q - 1}$;
 \item the elements $\alpha_p$ and $e_q$ satisfy
 \begin{equation*}
 \alpha_p e_q = \frac1{c_0}
 \left( \sum_{i = 1}^{n - 1}
 \frac{c_i \eps^{i(q - p)} g^i}{\eps^i - 1}
 \right) e_q.
 \end{equation*}
 \end{itemize}
\end{Lemma}

\begin{Lemma}
 \label{lem:xy_poly}
 For all positive integers $a$,
 the elements $x^a y^a$ and $y^a x^a$
 are polynomials in $k$ and group elements
 given by
 \begin{align*}
 x^a y^a ={}& \prod_{i = 0}^{a - 1} \Bigg( c_0 (k + i)
 + \sum\limits_{q = 0}^{n - 1} \beta_{q + i} e_q \Bigg),
 \\
 y^a x^a ={}& (-1)^a \prod\limits_{i = 0}^{a - 1}
 \Bigg( c_0 (-k + i)
 + \sum\limits_{q = 0}^{n - 1} \beta_{q + i} e_{-q}
 \Bigg).
 \end{align*}
 Also, the elements $x^a y^a e_q$ and $y^a x^a e_q$
 can be written as polynomial expressions in $k$
 multiplied by $e_q$ as follows
 \begin{align*}
 x^a y^a e_q ={}& e_q \prod\limits_{i = 0}^{a - 1}
 (c_0 (k + i) + \beta_{q + i}), \\
 y^a x^a e_q ={}& (-1)^a e_q \prod\limits_{i = 0}^{a - 1}
 (c_0 (-k + i) + \beta_{-q + i}).
 \end{align*}
\end{Lemma}

\begin{Lemma}
 \label{lem:trace_eh_zero}
 If $m \nmid q$, then $T(e_q h) = 0$.
\end{Lemma}

\begin{Theorem}
 \label{thm:classif}
 A linear map $T\colon H_c \to \CC$
 is a trace on $H_c$ if and only if
 all of the following conditions hold:
 \begin{enumerate}[label=\numbers]\itemsep=0pt
 \item $T$ is invariant under conjugation
 by elements of $\DD_n$;
 \label{cond:inv}
 \item $T\bigl(x^i R(k) e_q\bigr) = T\bigl(y^i R(k) e_q\bigr) = 0$
 when $i \neq 0$;
 \label{cond:g0}
 \item $T\bigl(x^i R(k) e_q h\bigr) = T\bigl(y^i R(k) e_q h\bigr) = 0$
 when $2 \nmid i$;
 \label{cond:h0}
 \item $T\bigl(x^{2i} R(k) e_q h\bigr) = T(S(k) R(k + i) e_{q - i} h)$,
 where $S$ is the polynomial
 satisfying that $x^i y^i e_{q - i}\allowbreak = S(k) e_{q - i}$;
 \label{cond:h_av}
 \item $T\bigl(y^{2i} R(k) e_q h\bigr) = (-1)^i T(S(k) R(k - i) e_{q + i} h)$,
 where $S$ is the polynomial
 satisfying that $y^i x^i e_{q + i} = S(k) e_{q + i}$;
 \label{cond:h_av_flip}
 \item $T\bigl(k^i e_q h\bigr) = 0$ for $i > 0$;
 \label{cond:k_poly}
 \item $T\bigl( \bigl(k + \half - \alpha_q\bigr) R\bigl(k + \half\bigr) e_q\bigr) =
 T\bigl( \bigl(k - \half - \alpha_q\bigr) R\bigl(k - \half\bigr) e_{q + 1}\bigr)$
 for all $q \in \ZZ$ and $R \in \CC[X]$;
 \label{cond:shift}
 \item $T( R(k) e_q ) = T( R(-k) e_{-q} )$
 for all $q \in \ZZ$ and $R \in \CC[X]$.
 \label{cond:even}
 \end{enumerate}
\end{Theorem}

\begin{proof}
 Assume first that $T$ is a trace;
 we will show that the listed conditions must hold.
 \begin{enumerate}[label=\numbers]\itemsep=0pt
 \item This is part of the definition of a trace.
 \item The trace condition gives
 $0 = T\bigl(\big[k, x^i R(k) e_q\big]\bigr) = -i T\bigl(x^i R(k) e_q\bigr)$,
 and similarly
 $0 = T\bigl(\big[k, y^i R(k) e_q\big]\bigr) = i T\bigl(y^i R(k) e_q\bigr)$,
 from which the claim follows.
 \item Conjugation by $h^2$ multiplies both $x$ and $y$ by $-1$
 while fixing $e_q$ and $h$, so
 \begin{gather*}
 T\bigl(x^i R(k) e_q h\bigr) = (-1)^i T\bigl(x^i R(k) e_q h\bigr) \qquad \text{and}\\
T\bigl(y^i R(k) e_q h\bigr) = (-1)^i T\bigl(y^i R(k) e_q h\bigr),\end{gather*}
 proving the claim.
 \item Using the trace condition
 to commute $x^i$ and $x^i R(k) e_q h$,
 we obtain
 \begin{align*}
 T\bigl(x^{2i} R(k) e_q h\bigr)
 ={}& T\bigl(x^i R(k) e_q h x^i\bigr)= T\bigl(x^i R(k) e_q y^i h\bigr) \\
 ={}& T\bigl(x^i y^i R(k + i) e_{q - i} h\bigr) = T\bigl(x^i y^i e_{q - i} R(k + i) h\bigr) \\
 ={}& T(S(k) e_{q - i} R(k + i) h) = T(S(k) R(k + i) e_{q - i} h),
 \end{align*}
 as desired.
 \item Similarly to the proof of condition~\ref{cond:h_av},
 we obtain
 \begin{align*}
 T\bigl(y^{2i} R(k) e_q h\bigr)
 ={}& T\bigl(y^i R(k) e_q h y^i\bigr)= (-1)^i T\bigl(y^i R(k) e_q x^i h\bigr) \\
 ={}& (-1)^i T\bigl(y^i x^i R(k - i) e_{q + i} h\bigr) = (-1)^i T(S(k) R(k - i) e_{q + i} h),
 \end{align*}
 as needed.
 \item The trace condition gives
 $0 = T\bigl(\big[k, k^{i - 1} e_q h\big]\bigr) = T\bigl(2 k^i e_q h\bigr)$,
 from which the claim follows.
 \item Using Lemma~\ref{lem:a_prop}
 and the commutativity of $e_q$ and $k$, we obtain
 \begin{align*}
 T\bigl( \bigl( k + \thalf - \alpha_q \bigr)
 R\bigl( k + \thalf \bigr) e_q \bigr)
 ={}& T\left( \left( k + \half -
 \frac1{c_0} \sum\limits_{i = 1}^{n - 1}
 \frac{c_i g^i}{\eps^i - 1} \right)
 R\bigl( k + \thalf \bigr) e_q \right) \\
 ={}& \frac1{c_0}
 T\bigl( xy R\bigl( k + \thalf \bigr) e_q \bigr).
 \end{align*}
 Using the trace condition
 to commute $x$ and $y R\bigl(k + \half\bigr) e_q$,
 we find
 \begin{align*}
 T\bigl( \bigl( k + \thalf - \alpha_q \bigr)
 R\bigl( k + \thalf \bigr) e_q \bigr)
 ={}& \frac1{c_0}
 T\bigl( y R\bigl( k + \thalf \bigr) e_q x \bigr)
 = \frac1{c_0}
 T\bigl( yx R\bigl( k - \thalf \bigr)
 e_{q + 1} \bigr) \\
 ={}& T\left( \left( k - \half -
 \frac1{c_0} \sum\limits_{i = 1}^{n - 1}
 \frac{c_i \eps^i g^i}{\eps^i - 1} \right)
 R\bigl( k - \thalf \bigr) e_{q + 1} \right),
 \end{align*}
 and using Lemma~\ref{lem:a_prop}
 and the commutativity of $e_q$ and $k$
 once again, we obtain
 \begin{align*}
 T\bigl( \bigl( k + \thalf - \alpha_q \bigr)
 R\bigl( k + \thalf \bigr) e_q \bigr)
 ={}& T\bigl( \bigl( k - \thalf - \alpha_q \bigr)
 R\bigl( k - \thalf \bigr) e_{q + 1} \bigr),
 \end{align*}
 as desired.
 \item This follows from conjugation by $h$.
 \end{enumerate}

 We now prove the converse.
 Suppose $T$ is a linear map from $H_c$ to $\CC$
 satisfying the listed conditions.
 We wish to show that for all $a$ in $H_c$,
 we have
 \[
 T([x, a]) = T([y, a]) = T([g, a]) = T([h, a]) = 0.
 \]
 The equality $T([g, a]) = T([h, a]) = 0$ follows
 from condition~\ref{cond:inv},
 so it remains to verify the result for commutation by $x$ and $y$.
 Also, the statements that $T([y, a]) = 0$ for all $a$
 follows from the statement that $T([x, a]) = 0$ for all $a$
 by conjugation by $h$ and condition~\ref{cond:inv}.
 It therefore suffices to show that $T([x, a]) = 0$
 when $a$ is equal to a basis element
 of the form~$x^i R(k) e_q$, $y^i R(k) e_q$,
 $x^i R(k) e_q h$, or $y^i R(k) e_q h$.
 We now check all possible cases.
 \begin{itemize}\itemsep=0pt
 \item Let $a = x^i R(k) e_q$ with $i$ nonnegative
 or $a = y^i R(k) e_q$ with $i > 1$.
 Then $T(xa) = 0 = T(ax)$ by
 condition~\ref{cond:g0}.
 \item Let $a = y R(k) e_q$.
 Then, using Lemma~\ref{lem:a_prop}, we obtain
 \begin{align*}
 T(xa) ={}& T\left( xy R(k) e_q \right) = c_0 T\left( \left( k + \half -
 \frac1{c_0} \sum\limits_{i = 1}^{n - 1}
 \frac{c_i g^i}{\eps^i - 1} \right)
 R(k) e_q \right) \\
 ={}& c_0 T\bigl( \bigl( k + \thalf
 - \alpha_q \bigr) R(k) e_q \bigr).
 \end{align*}
 Using condition~\ref{cond:shift}
 and Lemma~\ref{lem:a_prop},
 we obtain
 \begin{align*}
 T(xa) ={}& c_0
 T\bigl( \bigl( k - \thalf - \alpha_q \bigr)
 R(k - 1) e_{q + 1} \bigr) \\
 ={}& c_0 T\left( \left( k - \half -
 \frac1{c_0} \sum\limits_{i = 1}^{n - 1}
 \frac{c_i \eps^i g^i}{\eps^i - 1} \right)
 R(k - 1) e_{q + 1} \right) \\
 ={}& T\left( yx R(k - 1) e_{q + 1} \right)= T\left( y R(k) e_q x \right)= T(ax),
 \end{align*}
 as desired.
 \item Let $a = x^i R(k) e_q h$
 or $a = y^i R(k) e_q h$,
 with $i$ even.
 Then by condition~\ref{cond:h0},
 $T(xa) = 0 = T(ax)$.
 \item Let $a = x^{2i - 1} R(k) e_q h$,
 and let $S(X)$ be the polynomial
 such that $S(k) e_{q - 1} = xy e_{q - 1}$.
 Then by condition~\ref{cond:h_av},
 \begin{align*}
 T(xa) ={}& T\bigl( x^{2i} R(k) e_q h \bigr)= T\bigl( x^i y^i R(k + i)
 e_{q - i} h \bigr) \\
 ={}& T\bigl( x^{i - 1} (xy) y^{i - 1} R(k + i)
 e_{q - i} h \bigr) = T\bigl( x^{i - 1} (xy) y^{i - 1}
 e_{q - i} R(k + i) h \bigr) \\
 ={}& T\bigl( x^{i - 1} (xy) e_{q - 1} y^{i - 1}
 R(k + i) h\bigr) = T\bigl( x^{i - 1} S(k) e_{q - 1} y^{i - 1}
 R(k + i) h \bigr) \\
 ={}& T\bigl( x^{i - 1} y^{i - 1}
 S(k + i - 1) R(k + i) e_{q - i} h\bigr) = T\bigl( x^{2i - 2} S(k)
 R(k + 1) e_{q - 1} h \bigr) \\
 ={}& T\bigl( x^{2i - 2} (xy)
 R(k + 1) e_{q - 1} h \bigr) = T\bigl( x^{2i - 1}
 R(k) e_q y h \bigr) \\
 ={}& T\bigl(x^{2i - 1} R(k) e_q h x\bigr) = T(ax),
 \end{align*}
 as needed.
 \item Let $a = y^{2i + 1} R(k) e_q h$,
 and let $S(X)$ be the polynomial
 such that $S(k) e_{q + 2i} = xy e_{q + 2i}$.
 Then by condition~\ref{cond:h_av_flip},
 \begin{align*}
 T(xa) ={}& T\bigl( x y^{2i + 1}
 R(k) e_q h \bigr) = T\bigl( xy e_{q + 2i} y^{2i}
 R(k) h \bigr) \\
 ={}& T\bigl( S(k) e_{q + 2i} y^{2i}
 R(k) h \bigr) = T\bigl( y^{2i}
 S(k + 2i) R(k) e_q h \bigr) \\
 ={}& (-1)^i T\bigl( y^i x^i
 S(k + i) R(k - i) e_{q + i} h \bigr).
 \end{align*}
 On the other hand,
 \begin{align}
 T(ax) ={}& T\bigl( y^{2i + 1}
 R(k) e_q y h \bigr) = T\bigl( y^{2i + 2}
 R(k + 1) e_{q - 1} h \bigr) \nonumber\\
 ={}& (-1)^{i + 1} T\bigl( y^{i + 1} x^{i + 1}
 R(k - i) e_{q + i} h \bigr).\label{eq:t_ax}
 \end{align}
 So by condition~\ref{cond:k_poly},
 both $T(xa)$ and $T(ax)$
 are proportional to $T(e_{q + i} h)$,
 which, by Lemma~\ref{lem:trace_eh_zero},
 is nonzero only if $m \mid q + i$.
 We may therefore assume that $m \mid q + i$,
 because otherwise, $T(xa) = 0 = T(ax)$.
 It follows that $n \mid 2q + 2i$,
 so $\beta_{q + 2i} = \beta_{-q}$.
 Therefore, by Lemma~\ref{lem:xy_poly},
 \begin{align*}
 T(xa)
 ={}& (-1)^i T\bigl( y^i x^i
 S(k + i) R(k - i) e_{q + i} h \bigr) \\\
 ={}& T\Bigg( \Bigg(
 \prod\limits_{j = 0}^{i - 1}
 c_0 (-k + j) + \beta_{-q - i + j}
 \Bigg)
( c_0 (k + i) + \beta_{q + 2i})
 R(k - i) e_{q + i} h \Bigg) \\
 ={}& T\Bigg( \Bigg(
 \prod\limits_{j = 0}^{i - 1}
 c_0 (-k + j) + \beta_{-q - i + j}
 \Bigg)
( c_0 (k + i) + \beta_{-q})
 R(k - i) e_{q + i} h \Bigg).
 \end{align*}
 By condition~\ref{cond:k_poly},
 this expression depends only on its value
 for $k = 0$.
 Therefore, simplifying the expression for $T(xa)$
 and using \eqref{eq:t_ax}, we obtain
 \begin{align*}
 T(xa)
 ={}& T\Bigg( \Bigg(
 \prod\limits_{j = 0}^{i - 1}
 c_0 (-k + j) + \beta_{-q - i + j}
 \Bigg)
( c_0 (-k + i) + \beta_{-q})
 R(k - i) e_{q + i} h \Bigg) \\
 ={}& T\Bigg( \Bigg(
 \prod\limits_{j = 0}^i
 c_0 (-k + j) + \beta_{-q - i + j}
 \Bigg)
 R(k - i) e_{q + i} h \Bigg) \\
 ={}& (-1)^{i + 1} T\bigl( y^{i + 1} x^{i + 1}
 R(k - i) e_{q + i} h \bigr) = T(ax),
 \end{align*}
 as needed.
 \end{itemize}
 This resolves all cases, completing the proof.
\end{proof}

\end{document}
